% ============================================================
% Background / Preliminaries
% Pedagogical introduction to the concepts the thesis relies on.
% Assumes no prior familiarity with the specific task. Uses plain
% \cite{}; all keys resolve against mainbib.bib.
% ============================================================

\section{Background}
\label{sec:background}

This section introduces the concepts and terminology used throughout the thesis. It assumes general
machine-learning knowledge but not prior familiarity with LLMs, fact-checking, ranking, multilingual
evaluation, or the metrics used here. We first explain language models and their adaptation to a task.
We then introduce text ranking, binary classification, numerical fact-checking, cross-lingual
consistency, and numerical reasoning. The section ends with the task and agreement metrics used in
the experiments.

\subsection{Large Language Models}
\label{subsec:bg-llm}

A \emph{large language model} (LLM) is a neural network that models the probability of natural-language
text. Modern LLMs are usually decoder-only Transformers~\cite{vaswani2017attention}. They learn by
\emph{next-token prediction}: given the preceding tokens, the model predicts the next one. For a
sequence $x_1,\dots,x_T$, the joint probability is
$\prod_{t=1}^{T} p_\theta(x_t \mid x_{<t})$. During \emph{pre-training}, the model maximises this
likelihood over a large text corpus. The resulting linguistic and world knowledge can then be adapted
to specific tasks.

\paragraph{Tokens and tokenisation.}
LLMs process \emph{tokens}, which are subword units produced by a \emph{tokeniser}. The tokeniser is
learned from the pre-training corpus. Common words may form one token, while rare words are split into
several pieces. This scheme, often based on byte-pair encoding~\cite{sennrich2016bpe}, keeps the
vocabulary finite while representing any string. Tokenisation matters here because numbers and
non-Latin scripts may be split differently, which can affect numerical reasoning across languages
(Section~\ref{subsec:bg-numeracy}).

\paragraph{Instruction tuning.}
A raw pre-trained model predicts plausible continuations but may not follow instructions reliably.
\emph{Instruction tuning} fine-tunes it on instruction--response examples. The resulting models can
follow natural-language requests. The three instruction-tuned \emph{backbones} used here are
Mistral-7B-Instruct~\cite{jiang2023mistral7b}, Llama-3.1-8B-Instruct~\cite{llama31instruct}, and
Qwen3-8B~\cite{qwen3technicalreport}. They can rank items, classify text, or answer questions without
task-specific training. The numbers 7B and 8B indicate billions of parameters, a rough measure of
model capacity.

\subsection{Adapting Language Models to a Task}
\label{subsec:bg-adapt}

An LLM can be adapted to a task in two broad ways. \emph{Prompting} leaves its weights unchanged,
whereas \emph{fine-tuning} updates some or all of them. This thesis uses both approaches.

\paragraph{Prompting and in-context learning.}
In \emph{zero-shot} prompting, the model receives only an instruction and the input. \emph{Few-shot}
prompting also provides several solved examples, allowing the model to infer the pattern through
\emph{in-context learning}. Prompting requires no gradient updates and is therefore a natural
baseline. However, its output can be sensitive to wording and formatting.

\paragraph{Chain-of-thought reasoning.}
\emph{Chain-of-thought} (CoT) prompting asks the model to produce intermediate reasoning steps before
its final answer~\cite{WeiWSB23}. This often improves multi-step tasks such as arithmetic and
multi-hop questions, while making the reasoning inspectable. The ``reasoning traces'' in this thesis
are such chains of thought: each is a step-by-step argument over a claim and its evidence that ends in
a proposed verdict. The task provides these traces. This thesis studies how to select and rank them,
not how to generate them.

\paragraph{Full fine-tuning versus parameter-efficient fine-tuning.}
\emph{Fine-tuning} continues training on labelled task data. Updating every parameter of a 7--8B
model is expensive. \emph{Parameter-efficient fine-tuning} (PEFT) instead freezes the pre-trained
weights and trains a small set of new parameters. Low-Rank Adaptation
(LoRA)~\cite{hu2021loralowrankadaptationlarge} learns two small low-rank matrices for each adapted
weight matrix. Their product is added to the frozen weight. Only these \emph{adapters} are trained.
The adapter rank $r$ controls the number of trainable parameters, and $\alpha$ scales the update.

\paragraph{Quantisation and QLoRA.}
\emph{Quantisation} stores weights at lower precision, such as 4~bits instead of 16 or 32, to reduce
memory use. QLoRA~\cite{dettmers2023qloraefficientfinetuningquantized} loads the frozen base model at
4-bit precision and trains higher-precision LoRA adapters on top. This makes billion-parameter
fine-tuning feasible within the available GPU budget. Figure~\ref{fig:lora-qlora} illustrates the
setup.

\begin{figure}[htbp]
  \centering
  \includesvg[width=\linewidth,height=0.4\textheight,keepaspectratio]{lora-qlora}
  \caption{LoRA and QLoRA. The pretrained weight matrix $W$ is frozen a small pair of low-rank
  matrices $A$ and $B$ (with rank $r\ll d$) is trained and their product is added to the frozen
  weight, so the layer computes $h = Wx + (\alpha/r)\,BAx$ and only $A$ and $B$ are learned. QLoRA
  additionally stores the frozen weight $W$ in 4-bit precision to reduce memory.}
  \label{fig:lora-qlora}
\end{figure}

\paragraph{Preference optimisation.}
Models can also learn from \emph{preferences}, a preferred and a rejected output for the same input.
Reinforcement Learning from Human Feedback (RLHF) requires a separate reward model and reinforcement
learning. Direct Preference Optimization (DPO) provides a simpler alternative
\cite{rafailov2024directpreferenceoptimizationlanguage}. It increases the likelihood of preferred
outputs relative to rejected ones while keeping the model close to a reference model. This thesis
uses DPO as a ranking baseline.

\subsection{Text Ranking with Language Models}
\label{subsec:bg-ranking}

\emph{Ranking} orders candidate items by how well they satisfy a criterion. \emph{Re-ranking} applies
this process to an already retrieved candidate set. LLM rankers use three common formulations, which
differ in how many candidates the model sees at once (Figure~\ref{fig:ranking-formulations}).

\begin{itemize}
\item \textbf{Pointwise.} Each candidate is scored independently, and the ranking is obtained by
sorting the scores. This is the formulation our trace scorer uses.
\item \textbf{Pairwise.} The model compares two candidates at a time and states which is better, a
global ordering is assembled from the accumulated pairwise
preferences~\cite{qin2024largelanguagemodelseffective}. This keeps each decision simple but requires a
number of comparisons that grows quadratically with the number of candidates.
\item \textbf{Listwise.} All candidates are presented in a single prompt and the model emits a complete
ordering. This is efficient at inference but couples the decision to the whole input, including the
order in which candidates are listed~\cite{liu-etal-2025-lipo}.
\end{itemize}

\begin{figure}[htbp]
  \centering
  \includesvg[width=\linewidth,height=0.42\textheight,keepaspectratio]{ranking-formulations}
  \caption{The three standard formulations for ranking candidates with an LLM. \emph{Pointwise}
  scoring evaluates each candidate independently and sorts by score (the formulation used by our
  trace scorer), \emph{pairwise} ranking compares two candidates at a time and orders them by
  accumulated wins, \emph{listwise} ranking presents all candidates in one prompt and reads off the
  model's output order.}
  \label{fig:ranking-formulations}
\end{figure}

\paragraph{LLM-as-a-judge and presentation bias.}
Using an LLM to evaluate candidates is often called \emph{LLM-as-a-judge}. Such judgments can depend
on factors unrelated to content quality. \emph{Position bias} favours a candidate because of where it
appears in the prompt. Related effects include verbosity bias and sensitivity to prompt
format~\cite{shi2025judgingjudgessystematicstudy,qiao2026llmbasedlistwisererankingeffect}. These
biases motivate the move from generative ranking to independently scored candidates.

\subsection{Classification and Training Objectives}
\label{subsec:bg-classification}

Our system reformulates ranking as \emph{binary classification}. A classifier maps an input to one or
more real-valued \emph{logits}. Binary classification can use one logit $\ell$ with the sigmoid
function $\sigma(\ell) = 1/(1+e^{-\ell})$, or two logits
$(\ell^{\text{neg}}, \ell^{\text{pos}})$ with softmax. Both produce a probability for the positive
class. Training minimises cross-entropy between this probability and the true label, the one-logit
case uses binary cross-entropy. The positive-class probability becomes the relevance score used for
ranking. A small \emph{classification head} converts the model's internal representation into the
required logits.

\subsection{Automated Fact-Checking and Numerical Claim Verification}
\label{subsec:bg-factcheck}

\emph{Automated fact-checking} assesses whether evidence supports a claim. A full pipeline usually
includes \emph{claim detection}, \emph{evidence retrieval}, and \emph{claim verification}. The first
step identifies checkworthy claims, the second finds relevant documents, and the third assigns a
veracity label from the evidence. This thesis assumes that claim detection and retrieval are complete
and focuses on verification~\cite{sahitaj2025automatedfactcheckingrealworldclaims}.

\paragraph{Core components: claims, evidence, and verdicts.}
A verification example contains a \emph{claim}, retrieved \emph{evidence}, and a \emph{verdict}. The
verdict describes the relationship between claim and evidence and is one of \textsf{True},
\textsf{False}, or \textsf{Conflicting}. The task adds a fourth component: candidate reasoning traces,
each with its own proposed verdict, which the system must rank. A \emph{numerical claim} depends on
quantities such as counts, dates, percentages, rankings, or comparisons. Small changes in a value,
unit, or comparison scope can reverse its verdict~\cite{v2024quantemprealworldopendomainbenchmark}.

\paragraph{Reasoning traces and verifier-guided verification.}
\emph{Test-time scaling} generates several independent reasoning traces instead of one answer. A
separate \emph{verifier} selects which traces to trust
\cite{chungkham2025thinkrightmoretesttime}. The task studied here follows this verifier-guided setup.
It provides the claim, evidence, traces, and per-trace verdicts. The system ranks the traces and then
derives a claim-level verdict from the highest-ranked ones. Because trace generation is outside the
system, this is a ranking-and-aggregation task rather than free-form reasoning.
Figure~\ref{fig:factcheck-pipeline} shows the complete pipeline.

\begin{figure}[htbp]
  \centering
  \includesvg[width=\linewidth,height=0.5\textheight,keepaspectratio]{factcheck-pipeline}
  \caption{Numerical claim verification framed as verifier-guided trace ranking. A numerical claim and
  its retrieved evidence, together with a set of candidate reasoning traces (each ending in a proposed
  verdict), are provided in the task. The system this thesis develops scores and ranks the traces,
  aggregates the top-ranked verdicts by majority vote, and outputs a final claim verdict from
  \{\textsf{True}, \textsf{False}, \textsf{Conflicting}\}.}
  \label{fig:factcheck-pipeline}
\end{figure}

\paragraph{Benchmarks and the shared task.}
QuanTemp provides real-world numerical claims and evidence for open-domain
verification~\cite{v2024quantemprealworldopendomainbenchmark}. CLEF~CheckThat!~2026 Task~2 extends the
verifier-guided formulation to English, Spanish, and Arabic~\cite{clef-checkthat:2026:task2}. Our
system was developed in this setting~\cite{khawaja2026gipplab}. Other multilingual benchmarks,
including X-Fact~\cite{gupta2021xfactnewbenchmarkdataset} and
XFEVER~\cite{chang2023xfeverexploringfactverification}, likewise show that verification transfers
poorly across languages and that performance varies widely.

\subsection{Multilinguality and Cross-Lingual Consistency}
\label{subsec:bg-multilingual}

A \emph{multilingual} model is trained on several languages. The aim is \emph{cross-lingual transfer},
an ability learned in one language should carry over to another. Multilingual models learn partly
shared representations, so equivalent inputs can be encoded similarly. Interpretability studies of
decoder-only LLMs suggest that non-English input passes through intermediate representations close to
English before returning to the target language. This is described as a \emph{latent
language}~\cite{wendler2024llamasworkenglishlatent,schut2025multilingualllmsthinkenglish}. Such shared
structure makes language-invariant behaviour possible in principle.

\paragraph{Consistency and language invariance.}
In practice, models are often inconsistent. \emph{Consistency} means that meaning-preserving input
changes produce non-contradictory outputs
\cite{elazar2021measuringimprovingconsistencypretrained}. When the change is a translation, this becomes
\emph{cross-lingual consistency}, or \emph{language invariance}
\cite{Fierro_2022,Qi_2023}. Consistency and accuracy are separate, a model can be consistently wrong,
or accurate on average while disagreeing across translations
\cite{elazar2021measuringimprovingconsistencypretrained,Qi_2023}. Matching answers also do not prove
that knowledge is represented identically inside the model
\cite{ifergan2024beneathsurfaceconsistencyexploring}. The evaluation must therefore measure invariance
and accuracy separately.

\subsection{Numerical Reasoning in Language Models}
\label{subsec:bg-numeracy}

\emph{Numeracy} is the ability to reason about numbers, and it remains a weak point of language
models. Models often fail on elementary operations
\cite{rahman2025fragilenumbersenseprobing,wallace2019nlpmodelsknownumbers}. One reason is that standard
token embeddings do not explicitly encode numerical magnitude. Subword tokenisation can also split a
numeral into pieces that obscure its value~\cite{wallace2019nlpmodelsknownumbers}. \emph{Value-aware}
representations address this by encoding the numeric value of a span
\cite{dutulescu2026valueawarenumericalrepresentationstransformer,schwartz2024numerologicnumberencodingenhanced}.
Script and formatting create a separate problem. Arabic-Indic digits and different decimal or
thousands separators lead to different tokenisations, and accuracy can fall even when the arithmetic
is unchanged~\cite{reddy2026effectscriptsformatsllm,bhattacharya2025investigatinginteractionlinguisticmathematical}.
This makes numeral script an important source of variation for Arabic.

\subsection{Evaluation Metrics and Agreement Statistics}
\label{subsec:bg-metrics}

Finally, we define the metrics used to report results. They fall into three groups, classification
quality, ranking quality, and cross-lingual agreement.

\paragraph{Classification metrics.}
For one class, \emph{precision} is the fraction of predicted members that are correct, while
\emph{recall} is the fraction of true members that are retrieved. Their harmonic mean is
$\mathrm{F1} = 2\,PR/(P+R)$. \emph{Macro-F1} averages F1 across classes and therefore gives each
class equal weight. We use it for the three-class claim verdict.

\paragraph{Ranking metrics.}
For a ranked list, \emph{Recall@$k$} is the fraction of relevant items in the top $k$, while
\emph{Precision@$k$} is the fraction of the top $k$ items that are relevant. \emph{Mean Reciprocal
Rank at $k$} (MRR@$k$) uses the reciprocal rank of the first relevant item, or zero if none appears in
the top $k$, and averages it over queries. A trace is relevant here when its verdict matches the gold
claim label. Recall@5 is the primary ranking metric.

\paragraph{Inter-rater agreement.}
We treat each language-conditioned verdict as a ``rater'' of the same claim. \emph{Raw agreement} is
the fraction of claims with matching labels, but it includes agreement expected by chance.
Chance-corrected statistics address this. Cohen's $\kappa$ compares two
raters~\cite{cohen1960coefficient}, Fleiss' $\kappa$ handles more than
two~\cite{fleiss1971measuring}, and Krippendorff's $\alpha$ can also tolerate missing
data~\cite{krippendorff2004content}. Each equals $1$ for perfect agreement and is around $0$ at chance
level.

\paragraph{Rank-agreement.}
Rankings require rank-correlation measures. Kendall's $\tau$ compares the number of concordant and
discordant item pairs~\cite{kendall1938new}, while Spearman's $\rho$ correlates the assigned
ranks~\cite{spearman1904proof}. Both range from $-1$ for reversed order through $0$ for unrelated
order to $+1$ for identical order. Jaccard overlap compares the two top-$k$ sets by dividing the size
of their intersection by the size of their union. These measures show whether a model ranks the same
traces similarly across languages, not only whether it reaches the same verdict.
Section~\ref{sec:methodology} uses these concepts to define the scorer, invariance objective, and
evaluation protocol.
